\section{Conditions and metrics}
Every condition answers the same question at step $t$: what text should the agent type? \textbf{Full history, prompted} sees the OCR text of every earlier screen and names the screen and the text in one call of the 14B reader. Its context is 1,241 tokens on average in the short set and 4,584 for the longest GUI-Odyssey trajectory. \textbf{Screen index, prompted} sees one 32-token line per screen written by Qwen2.5-7B-Instruct, names a screen, then reads it in full; it is ReadAgent's gist-and-lookup~\citep{readagent} on screens. \textbf{Per-screen scorer, prompted} runs our scorer's prompt on the 7B model without the adapter. \textbf{Full history, trained} and \textbf{per-screen scorer, trained} are the two instantiations of \S\ref{sec:method}. \S\ref{sec:index} adds the index's writer variants, a no-history condition, a per-screen summary, lexical retrieval (BM25), and an oracle over sampled index lines.

\emph{Localization} is correct when the named screen is any screen whose OCR text contains the gold text (\emph{any-valid scoring}); \emph{end-to-end} is correct when the produced text, normalized as in mining, contains the gold text. We pair all comparisons on the same cases and test them with McNemar's test, reporting the net number of cases that flip and the two-sided $p$.
